\section{Background and Related Work}

\subsection{Multi-Agent Software Engineering}

LLM-based multi-agent systems use role specialization, iterative review, and orchestration to address software-engineering tasks. He et al. identify communication quality, conflict resolution, misinformation control, and evaluation beyond final artifacts as central research challenges~\cite{he2025road}. MAST empirically organizes multi-agent failures into system-design failures, inter-agent misalignment, and task-verification failures~\cite{cemri2025mast}. SEMAP responds with explicit behavioral contracts, structured messaging, and lifecycle verification~\cite{mao2025semap}. These studies motivate our focus on the representation and verification of intermediate evidence rather than on introducing another role architecture.

Self-Refine is an important baseline because one model generates an answer, critiques it, and revises it without additional training~\cite{madaan2023selfrefine}. Comparing this workflow with role-specialized agents helps separate the effect of distinct roles from the effect of spending three model calls.

\subsection{LLM-Based Vulnerability Analysis}

SecVulEval provides real-world C/C++ vulnerable and fixed functions with statement-level change information and reports that current LLMs remain weak at vulnerability localization with correct reasoning~\cite{ahmed2025secvuleval}. GPTLens separates generation and discrimination to reduce false positives in smart-contract auditing~\cite{hu2023gptlens}. VulAgent uses hypothesis generation and validation for project-level vulnerability analysis~\cite{wang2026vulagent}. MulVul combines a retrieval-augmented router with specialized detector agents and cross-model prompt evolution~\cite{wu2026mulvul}. LAMPS demonstrates role-specialized multi-agent reasoning for malicious-package detection~\cite{zeshan2026lamps}. Collectively, this work establishes that critic and validator roles are common in agentic security analysis. Our contribution is not the first multi-agent detector; it is the controlled isolation of handoff format and the measurement of evidence propagation.

\subsection{Fault Injection and Resilience}

Huang et al. inject faulty agent messages and show that challenge and inspection mechanisms can improve resilience~\cite{huang2025resilience}. MAS-FIRE systematizes intra-agent and inter-agent fault types and evaluates how architectures recover from modified prompts, responses, and routing~\cite{jia2026masfire}. We adapt this idea to vulnerability evidence. The injected conditions target a verdict, the availability of supporting evidence, or a concrete false claim and location. Unlike end-task-only evaluation, our analysis records whether the injected content appears in the final report.

The closest conceptual expectation is that structure should reduce ambiguity and make verification easier. The competing expectation is that schema-compliant misinformation may be treated as authoritative. Our experiment tests both possibilities.
